% This is samplepaper.tex, a sample chapter demonstrating the
% LLNCS macro package for Springer Computer Science proceedings;
% Version 2.21 of 2022/01/12
%
\documentclass[runningheads]{llncs}
%
\usepackage{url}
\usepackage{marvosym}

\usepackage[printonlyused,withpage]{acronym}
\newacro{ViT}[ViT]{Vision Transformer}
\newacro{MHSA}[MHSA]{Multi-Head Self-Attention}
\newacro{LoRA}[LoRA]{Low-Rank Adaptation}
\newacro{QLoRA}[QLoRA]{Quantized Low-Rank Adaptation}
\newacro{CV}[CV]{computer vision}
\newacro{SOTA}[SOTA]{state-of-the-art}
\newacro{PEFT}[PEFT]{parameter-efficient fine-tuning}
\newacro{CNN}[CNN]{convolutional neural network}
\newacro{KD}[KD]{Knowledge Distillation}
\newacro{LLM}[LLM]{large language model}
\newacro{NLP}[NLP]{natural language processing}
\newacro{NF4}[NF4]{4-bit NormalFloat}
\newacro{DQ}[DQ]{Double Quantization}
\newacro{GELU}[GELU]{ Gaussian Error Linear Unit}
\newacro{FP32}[FP32]{32-bit floating-point}
\newacro{SVD}[SVD]{singular value decomposition}
\newacro{PTQ}[PTQ]{post-training quantization}
\newacro{QAT}[QAT]{quantization-aware training}

\usepackage[T1]{fontenc}
% T1 fonts will be used to generate the final print and online PDFs,
% so please use T1 fonts in your manuscript whenever possible.
% Other font encondings may result in incorrect characters.
%
\usepackage{graphicx}
% Used for displaying a sample figure. If possible, figure files should
% be included in EPS format.
%
% If you use the hyperref package, please uncomment the following two lines
% to display URLs in blue roman font according to Springer's eBook style:
%\usepackage{color}
%\renewcommand\UrlFont{\color{blue}\rmfamily}
%\urlstyle{rm}
%
\begin{document}
%
\title{LoftQ: ViT(temp title)}
%
%\titlerunning{Abbreviated paper title}
% If the paper title is too long for the running head, you can set
% an abbreviated paper title here
%
\author{Hikaru Fukusaka\inst{1}\and
Lin Meng\inst{2}$^{(\textrm{\Letter})}$}
%
\authorrunning{H. Fukusaka et al.}
% First names are abbreviated in the running head.
% If there are more than two authors, 'et al.' is used.
%
\institute{Graduate School of Science and Engineering, Ritsumeikan University, Kusatsu 5250058, Japan\and
College of Science and Engineering, Ritsumeikan University, Kusatsu 5250058, Japan\\
\email{menglin@fc.ritsumei.ac.jp}}
%
\maketitle              % typeset the header of the contribution
%
\begin{abstract}
The abstract should briefly summarize the contents of the paper in
150--250 words.

\keywords{ViT  \and PEFT \and LoRA \and Quantization.}
\end{abstract}
%
%
%
\section{Introduction}

\section{Related Work}
\subsection{Parameter-Efficient Fine-Tuning}
\ac{LoRA}~\cite{hu2022lora} freezes pre-trained weights and approximates updates $\Delta W \approx BA$ using low-rank matrices($r \ll d$).
The update is scaled by $\alpha/r$ to balance gradient magnitudes and ensure training stability,
drastically reducing trainable parameters while matching full fine-tuning performance.

AdaLoRA~\cite{zhang2023adaptive} overcomes \ac{LoRA}'s fixed-rank limitation by dynamically allocating ranks based on gradient importance.
It utilizes \ac{SVD} to prune negligible singular values and redistributes the freed ranks to weight matrices with larger gradient magnitudes, optimizing parameter efficiency without manual rank tuning.

%この章はちょっと微妙かも LoRA-FAの説明が若干違う
Several variants further enhance efficiency:
LoRA-FA employs a hybrid low-rank and full-rank structure to capture complex updates;
DoRA decomposes updates into magnitude and direction, controlling weight norms to prevent performance degradation;
and VeRA represents updates via vector outer products to minimize the parameter footprint.

In \ac{ViT}, \ac{PEFT} is specialized via VPT (Jia et al., 2022), which injects visual prompts; SSF (Lian et al., 2022),
which fine-tunes spectral-spatial features; and AdaptFormer (Chen et al., 2022), which inserts bottleneck adapters into self-attention and FFN blocks.
Standardized on the VTAB-1kbenchmark, these methods demonstrate that PEFT achieves full fine-tuning parity in vision tasks while mitigating catastrophic forgetting.

\subsection{Quantization for LLMs and quantization-aware fine-tuning}
\label{sec:rw:llm_quant}
Weight-only post-training quantization is the dominant strategy for deploying \acp{LLM}.
GPTQ~\cite{frantar2023gptqaccurateposttrainingquantization} uses approximate second-order error compensation to bring OPT~\cite{zhang2022optopenpretrainedtransformer}
and BLOOM~\cite{NEURIPS2022_ce9e92e3} down to INT4 with minor perplexity loss.
AWQ~\cite{MLSYS2024_42a452cb} observes that weights are not equally important: protecting roughly $1\%$ of salient channels---identified by input-activation magnitude rather than weight magnitude---greatly reduces weight quantization error.
It then applies a per-channel scaling that shrinks the relative quantization error of those channels while keeping the scheme weight-only.
A complementary thread treats activation outliers as a structural obstacle to mixed weight--activation quantization:
LLM.int8()~\cite{NEURIPS2022_c3ba4962} shows that emergent outlier features in large transformers break naive INT8 calibration,
and SmoothQuant\cite{pmlr-v202-xiao23c} migrates magnitude between activations and weights so that both sides become quantizer-friendly.
All of these methods are applied offline as post-training quantization to a frozen model,
and are not designed to interact with downstream fine-tuning.

\subsection{Quantization for Vision Transformers}
\label{sec:rw:vit_quant}
Quantization research tailored to ViTs has developed along two largely separate lines.
The first targets \ac{PTQ} and focuses on the unusual activation distributions induced by transformer non-linearities.
PTQ4ViT~\cite{10.1007/978-3-031-19775-8_12} observes that post-Softmax values are tightly concentrated near zero within $[0,1]$ while post-GELU activations are asymmetric,
with an unbounded positive range and a narrow negative one, and proposes twin uniform quantization with a Hessian-guided metric to handle both cases.
RepQ-ViT~\cite{Li_2023_ICCV} addresses post-LayerNorm activations with severe inter-channel variation and post-Softmax activations exhibiting power-law features by applying channel-wise and $\log\sqrt{2}$ quantizers during calibration,
and then reparameterizing the scales into hardware-friendly layer-wise and $\log_2$ quantizers at inference.
These methods push INT4 \ac{PTQ} to a usable accuracy range without modifying the backbone,
but by construction they do not interact with downstream fine-tuning.

A second line pursues low-bit \acp{ViT} through \ac{QAT}.
Q-ViT~\cite{NEURIPS2022_deb921bf} identifies the information distortion of the quantized self-attention map as the principal bottleneck and restores it through an information rectification module that maximizes the entropy of quantized attention weights,
combined with a distribution-guided distillation loss that aligns attention similarity with a full-precision teacher.
Subsequent fully binarized variants such as BiViT~\cite{He_2023_ICCV} push this regime further.
While these works show that aggressively low-bit \acp{ViT} are trainable,
they require end-to-end re-optimization of the entire backbone and do not exploit the parameter-efficient adaptation paradigm that has become standard for vision transfer learning~\cite{10.1007/978-3-031-19827-4_41,NEURIPS2022_69e2f49a}.
Across both lines, the patch embedding and the classification head are conventionally excluded from quantization, reflecting an empirical sense that the input projection is disproportionately fragile---an observation we revisit quantitatively in Section~\ref{sec:analysis}.
None of these efforts examine whether a LoftQ-style coupling between weight quantization and LoRA initialization transfers to \acp{ViT} under sub-2-bit budgets.
\section{Method}

\section{Experiments}

\section{Analysis}

\section{Conclusion}

%
% the environments 'definition', 'lemma', 'proposition', 'corollary',
% 'remark', and 'example' are defined in the LLNCS documentclass as well.
%
\begin{credits}
\subsubsection{\ackname}
This work was supported by JST SPRING, Grant Number JPMJSP2101.
\end{credits}
%
% ---- Bibliography ----
%
% BibTeX users should specify bibliography style 'splncs04'.
% References will then be sorted and formatted in the correct style.
%
\bibliographystyle{splncs04}
\bibliography{references}
\end{document}
